% ============================================================
\chapter{Bayesian Prediction}
\label{ch:prediction}
% ============================================================

In frequentist statistics, prediction uses a point estimate: one estimates $\theta$, then predicts $\tilde{y}$ as if $\hat\theta$ were the true value. This approach ignores parameter uncertainty and produces prediction intervals that are too narrow. The Bayesian approach is fundamentally different: instead of fixing $\theta$, one \emph{integrates} over all possible values of $\theta$, weighted by the posterior distribution. The resulting predictive distribution naturally incorporates two sources of uncertainty---the intrinsic randomness of the data and the uncertainty about the model---and produces better-calibrated forecasts.

\begin{intuition}[title=\textbf{Central idea}]
Bayesian prediction integrates parameter uncertainty by averaging the
predictive distribution over the posterior. This produces well-calibrated
prediction intervals and avoids the overfitting inherent in using point
estimates.
\end{intuition}

% ------------------------------------------------------------
\section{Prior predictive distribution}
% ------------------------------------------------------------

\begin{definition}[Prior predictive distribution]
The \textbf{prior predictive} (or marginal) distribution of a new
observation $\tilde{x}$ is:
\[
  m(\tilde{x}) = \int_\Theta f(\tilde{x} \mid \theta)\,\pi(\theta)\,d\theta.
\]
It does not depend on the data and represents the prediction before any
observation.
\end{definition}

\begin{example}
\textbf{Bernoulli--Beta model.}
With $\theta \sim \text{Be}(a, b)$ and
$\tilde{X} \mid \theta \sim \text{Ber}(\theta)$:
\[
  m(\tilde{x} = 1) = \E[\theta] = \frac{a}{a + b}.
\]
For the Laplace prior $\text{Be}(1, 1)$: $m(\tilde{x} = 1) = 1/2$.
\end{example}

\begin{example}
\textbf{Normal--Normal model.}
With $\mu \sim \mathcal{N}(\mu_0, \tau_0^2)$ and
$\tilde{X} \mid \mu \sim \mathcal{N}(\mu, \sigma^2)$:
\[
  m(\tilde{x}) = \mathcal{N}(\tilde{x} \mid \mu_0, \sigma^2 + \tau_0^2).
\]
The predictive variance $\sigma^2 + \tau_0^2$ combines aleatoric and
epistemic uncertainty.
\end{example}

% ------------------------------------------------------------
\section{Posterior predictive distribution}
% ------------------------------------------------------------

\begin{definition}[Posterior predictive distribution]
The \textbf{posterior predictive distribution} of $\tilde{x}$ given
$x = (x_1, \ldots, x_n)$ is:
\[
  f(\tilde{x} \mid x) = \int_\Theta f(\tilde{x} \mid \theta)\,\pi(\theta \mid x)\,d\theta.
\]
\end{definition}

\begin{keyformulas}[title=\textbf{Predictive variance decomposition}]
\[
  \Var[\tilde{X} \mid x]
  = \underbrace{\E_{\theta \mid x}[\Var[\tilde{X} \mid \theta]]}_{\text{aleatoric uncertainty}}
  + \underbrace{\Var_{\theta \mid x}[\E[\tilde{X} \mid \theta]]}_{\text{epistemic uncertainty}}.
\]
This decomposition is the \textbf{law of total variance}.
\end{keyformulas}

\begin{theorem}[Normal--Normal posterior predictive]
If $X_i \overset{\text{iid}}{\sim} \mathcal{N}(\mu, \sigma^2)$ with
$\mu \mid x \sim \mathcal{N}(\mu_n, \tau_n^2)$, then:
\[
  \tilde{X} \mid x \sim \mathcal{N}(\mu_n, \sigma^2 + \tau_n^2).
\]
\end{theorem}

\begin{proof}
Conditionally on $\mu$, $\tilde{X} \sim \mathcal{N}(\mu, \sigma^2)$.
Integrating over $\mu \mid x \sim \mathcal{N}(\mu_n, \tau_n^2)$, the
sum of two independent Gaussians gives:
\[
  \tilde{X} \mid x \sim \mathcal{N}(\mu_n, \sigma^2 + \tau_n^2).
\]
\end{proof}

\begin{theorem}[Bernoulli--Beta posterior predictive]
If $X_i \overset{\text{iid}}{\sim} \text{Ber}(\theta)$ with
$\theta \mid x \sim \text{Be}(a_n, b_n)$, then:
\[
  \PP(\tilde{X} = 1 \mid x) = \E[\theta \mid x] = \frac{a_n}{a_n + b_n}.
\]
This is \textbf{Laplace's rule of succession} when $a = b = 1$:
$\PP(\tilde{X} = 1 \mid x) = (s+1)/(n+2)$.
\end{theorem}

\begin{theorem}[Poisson--Gamma posterior predictive]
If $X_i \overset{\text{iid}}{\sim} \text{Poi}(\lambda)$ with
$\lambda \mid x \sim \text{Ga}(\alpha_n, \beta_n)$, the predictive is
a \textbf{negative binomial}:
\[
  \tilde{X} \mid x \sim \text{NegBin}\!\left(\alpha_n, \frac{\beta_n}{\beta_n + 1}\right).
\]
\end{theorem}

\begin{proof}
\begin{align*}
  \PP(\tilde{X} = k \mid x)
  &= \int_0^\infty \frac{e^{-\lambda}\lambda^k}{k!}
  \cdot \frac{\beta_n^{\alpha_n}}{\Gamma(\alpha_n)}
  \lambda^{\alpha_n - 1} e^{-\beta_n \lambda}\,d\lambda \\
  &= \frac{\beta_n^{\alpha_n}}{k!\,\Gamma(\alpha_n)}
  \cdot \frac{\Gamma(k + \alpha_n)}{(\beta_n + 1)^{k + \alpha_n}} \\
  &= \binom{k + \alpha_n - 1}{k}
  \left(\frac{\beta_n}{\beta_n + 1}\right)^{\alpha_n}
  \left(\frac{1}{\beta_n + 1}\right)^k.
\end{align*}
\end{proof}

\begin{theorem}[Normal--NIG posterior predictive]
If $(\mu, \sigma^2) \mid x \sim \text{NIG}(\mu_n, \kappa_n, \alpha_n, \beta_n)$,
the posterior predictive is a Student-$t$:
\[
  \tilde{X} \mid x \sim t_{2\alpha_n}\!\left(\mu_n,\;
  \frac{\beta_n(\kappa_n + 1)}{\alpha_n \kappa_n}\right).
\]
\end{theorem}

% ------------------------------------------------------------
\section{Posterior predictive checking}
% ------------------------------------------------------------

\begin{definition}[Posterior predictive check]
A \textbf{posterior predictive check} (PPC) simulates replicate datasets
$x^{\text{rep}}$ from the predictive distribution and compares them to
the observed data to assess model adequacy.
\end{definition}

\begin{definition}[Bayesian predictive $p$-value]
The \textbf{predictive $p$-value} for a test statistic $T(x)$ is:
\[
  p_B = \PP(T(x^{\text{rep}}) \geq T(x_{\text{obs}}) \mid x_{\text{obs}}).
\]
Values near $0$ or $1$ suggest poor model fit.
\end{definition}

\begin{algorithme}[title=\textbf{Posterior predictive check algorithm}]
\begin{enumerate}
  \item For $t = 1, \ldots, T$:
  \begin{enumerate}
    \item Draw $\theta^{(t)} \sim \pi(\theta \mid x_{\text{obs}})$ (from MCMC).
    \item Draw $x^{\text{rep},(t)} \sim f(x \mid \theta^{(t)})$.
    \item Compute $T^{(t)} = T(x^{\text{rep},(t)})$.
  \end{enumerate}
  \item Estimate $p_B = \frac{1}{T}\sum_{t=1}^T \mathbf{1}(T^{(t)} \geq T(x_{\text{obs}}))$.
  \item Plot the distribution of $T^{(t)}$ and compare to $T(x_{\text{obs}})$.
\end{enumerate}
\end{algorithme}

% ------------------------------------------------------------
\section{Prediction intervals}
% ------------------------------------------------------------

\begin{definition}[Prediction interval]
A $100(1-\alpha)\%$ \textbf{prediction interval} for $\tilde{X}$ is an
interval $[a, b]$ such that:
\[
  \PP(\tilde{X} \in [a, b] \mid x) = 1 - \alpha.
\]
\end{definition}

\begin{remark}
The prediction interval is always wider than the credible interval for
$\theta$, because it integrates both parameter uncertainty and the
randomness of the new observation.
\end{remark}

% ------------------------------------------------------------
\section{Bayesian vs.\ plug-in prediction}
% ------------------------------------------------------------

\begin{warning}[title=\textbf{Plug-in danger}]
The plug-in approach replaces $\theta$ by a point estimate $\hat{\theta}$
and predicts via $f(\tilde{x} \mid \hat{\theta})$. It underestimates
predictive uncertainty because it ignores parameter uncertainty.
Bayesian prediction is better calibrated.
\end{warning}

% ------------------------------------------------------------
\section{Python implementation}
% ------------------------------------------------------------

\begin{codeblock}[title=\textbf{Posterior predictive distribution}]
\begin{minted}{python}
import numpy as np
import pymc as pm
import arviz as az
import matplotlib.pyplot as plt

# Gaussian data
np.random.seed(42)
n = 30
mu_true, sigma_true = 5.0, 2.0
data = np.random.normal(mu_true, sigma_true, n)

# Bayesian model
with pm.Model() as model_pred:
    mu = pm.Normal("mu", mu=0, sigma=10)
    sigma = pm.HalfNormal("sigma", sigma=10)
    y = pm.Normal("y", mu=mu, sigma=sigma, observed=data)
    # Predictive
    y_pred = pm.Normal("y_pred", mu=mu, sigma=sigma)
    trace = pm.sample(3000, random_seed=42)
    ppc = pm.sample_posterior_predictive(trace,
                                         random_seed=42)

# Visualization
fig, axes = plt.subplots(1, 2, figsize=(12, 4))

az.plot_ppc(ppc, observed_adata=trace, ax=axes[0])
axes[0].set_title("Posterior predictive check")

y_pred_samples = ppc.posterior_predictive["y_pred"].values.flatten()
ci_pred = np.percentile(y_pred_samples, [2.5, 97.5])
axes[1].hist(y_pred_samples, bins=50, density=True,
             alpha=0.5, label="Predictive")
axes[1].axvline(ci_pred[0], color="red", linestyle="--",
                label=f"95% PI: [{ci_pred[0]:.1f}, {ci_pred[1]:.1f}]")
axes[1].axvline(ci_pred[1], color="red", linestyle="--")
axes[1].legend()
axes[1].set_title("Posterior predictive distribution")

plt.tight_layout()
plt.savefig("predictive.pdf")
\end{minted}
\end{codeblock}

\begin{codeblock}[title=\textbf{Bayesian predictive $p$-value}]
\begin{minted}{python}
# Test statistic: standard deviation
T_obs = data.std()

# Simulate replicates
T_rep = []
mu_samples = trace.posterior["mu"].values.flatten()
sigma_samples = trace.posterior["sigma"].values.flatten()
for i in range(len(mu_samples)):
    x_rep = np.random.normal(mu_samples[i],
                              sigma_samples[i], n)
    T_rep.append(x_rep.std())
T_rep = np.array(T_rep)

# Predictive p-value
p_B = np.mean(T_rep >= T_obs)
print(f"T(x_obs) = {T_obs:.3f}")
print(f"Predictive p-value = {p_B:.3f}")

plt.figure(figsize=(8, 4))
plt.hist(T_rep, bins=50, density=True, alpha=0.5,
         label="T(x_rep)")
plt.axvline(T_obs, color="red",
            label=f"T(x_obs) = {T_obs:.2f}")
plt.xlabel("Standard deviation")
plt.ylabel("Density")
plt.legend()
plt.title(f"Predictive p-value = {p_B:.3f}")
plt.tight_layout()
plt.savefig("ppc_pvalue.pdf")
\end{minted}
\end{codeblock}

\begin{codeblock}[title=\textbf{Poisson--Gamma predictive (negative binomial)}]
\begin{minted}{python}
from scipy import stats

# Poisson data
np.random.seed(123)
data_pois = np.random.poisson(3.0, size=20)

# Gamma(2, 1) prior
alpha_n = 2 + data_pois.sum()
beta_n = 1 + len(data_pois)
p_nb = beta_n / (beta_n + 1)

# Predictive: NegBin(alpha_n, p_nb)
x_grid = np.arange(0, 15)
pred_pmf = stats.nbinom.pmf(x_grid, alpha_n, p_nb)

plt.figure(figsize=(8, 4))
plt.bar(x_grid, pred_pmf, alpha=0.7,
        label="NegBin predictive")
plt.xlabel(r"$\tilde{x}$")
plt.ylabel("Probability")
plt.title("Poisson-Gamma predictive")
plt.legend()
plt.tight_layout()
plt.savefig("pred_negbin.pdf")
\end{minted}
\end{codeblock}

% ------------------------------------------------------------
\section{Exercises}
% ------------------------------------------------------------

\begin{exercise}
Derive the posterior predictive distribution for the Bernoulli--Beta
model for $m$ new observations (not just one). Show that it is a
beta-binomial distribution.
\end{exercise}

\begin{exercise}
Show that the predictive variance is always greater than or equal to
the conditional variance $\Var[\tilde{X} \mid \theta]$ for any value of
$\theta$.
\end{exercise}

\begin{exercise}
Derive the predictive distribution for the Normal--NIG model and show
that it is a Student-$t$.
\end{exercise}

\begin{exercise}
\textbf{(Numerical)} Simulate Poisson data and compare:
\begin{enumerate}
  \item the Bayesian predictive (negative binomial),
  \item the plug-in predictive (Poisson with $\hat{\lambda} = \bar{x}$).
\end{enumerate}
Compute the actual coverage of $95\%$ prediction intervals for both
approaches over $1000$ repetitions.
\end{exercise}

\begin{exercise}
Perform a posterior predictive check for a Poisson model applied to
overdispersed data. Use the variance as the test statistic.
\end{exercise}
